\documentclass[12pt]{article}

\usepackage{sbc-template}
\usepackage{graphicx,url}
\usepackage[utf8]{inputenc}
\usepackage[spanish]{babel}
\addto\captionsspanish{\renewcommand{\tablename}{Tabla}}
\usepackage{booktabs}
\usepackage{tikz}
\usetikzlibrary{positioning,arrows.meta}

\sloppy

\title{Hacia la Evolución de la GUI Desktop:\\ IA Centrada en el Humano como Paradigma\\ para Repensar la Interacción con Sistemas Operativos}

\author{Jose Ramon Aragon Toledo\inst{1}, Ricardo Ruiz Rodríguez\inst{1}}

\address{Universidad Tecnológica de la Mixteca (UTM)\\
  Huajuapan de León, Oaxaca -- México
  \email{aatr010423@gmail.com, rruiz@mixteco.utm.mx}
}

\begin{document}

\maketitle

\begin{abstract}
Innovation in desktop operating system interfaces has stagnated over the past two
decades. Despite exponential growth in computational power, tools such as file
explorers, terminals, and configuration managers have scarcely evolved since the
1990s. This plateau particularly affects GNU/Linux, whose power remains inaccessible
to non-expert users due to command-line interaction models from the 1970s. This paper
analyzes desktop interaction stagnation through the lens of Human-Centered AI (HCAI)
and presents BOTAS, a voice-based assistant enabling Spanish-speaking users to
interact with GNU/Linux via natural language. Evaluation across 200 tasks shows
89\% command-generation accuracy and 100\% detection of destructive operations.
\end{abstract}

\begin{resumo}
La innovación en las interfaces de sistemas operativos de escritorio ha
experimentado un estancamiento en las últimas dos décadas. A pesar del crecimiento
exponencial en capacidad computacional, herramientas como exploradores de archivos,
terminales y gestores de configuración apenas han evolucionado desde los años
noventa. Esta meseta afecta especialmente a GNU/Linux, cuyo potencial permanece
inaccesible para usuarios no expertos debido a modelos de interacción de línea de
comandos de los años setenta. Este artículo analiza el estancamiento de la
interacción con el escritorio a través del marco de la Inteligencia Artificial
Centrada en el Humano (HCAI) y presenta BOTAS, un asistente de voz que permite a
usuarios hispanohablantes interactuar con GNU/Linux mediante lenguaje natural. La
evaluación sobre 200 tareas muestra un 89\% de precisión en la generación de
comandos y un 100\% de detección de operaciones destructivas.
\end{resumo}


\section{Introducción}
\label{sec:intro}

\noindent
Operamos con procesadores cien millones de veces más rápidos que los de la década
de 1980, y sin embargo la interacción fundamental con un sistema operativo de
escritorio apenas ha cambiado en veinte años \cite{Jenson2024}. Scott
Jenson---veterano diseñador de experiencia de usuario que participó en el
desarrollo del Finder original de Macintosh en Apple y posteriormente condujo
investigación sobre interacción en Google---sostiene que ``este inmovilismo no es
un límite técnico sino una claudicación del diseño'' \cite{Jenson2024}. La
industria ha transitado de la transformación creativa que definió los albores de
la computación---donde conceptos de Xerox PARC fueron radicalmente
reimaginados---a una imitación superficial entre plataformas.

GNU/Linux sustenta aproximadamente el 90\% de la infraestructura en la nube y el
100\% de las 500 supercomputadoras más potentes del mundo
\cite{W3Techs2023, Top5002023}; no obstante, su interfaz de línea de comandos
perpetúa un paradigma de la década de 1970 \cite{Raymond2003} que exige memorizar
sintaxis crípticas sin retroalimentación sobre el riesgo de las operaciones.
Mientras tanto, los asistentes de voz comerciales gestionan consultas triviales
pero ninguno aborda la interacción profunda con el sistema operativo
\cite{Luger2016}, y herramientas como ChatGPT sugieren comandos sin ejecutarlos
ni validar su seguridad. Este trabajo sostiene que la Inteligencia Artificial
Centrada en el Humano (HCAI) \cite{Shneiderman2020, Shneiderman2022} ofrece un
marco para superar este estancamiento, y presenta BOTAS, un asistente de voz
que materializa los principios HCAI para la interacción con GNU/Linux en español.


\section{Marco Conceptual y Trabajo Relacionado}
\label{sec:framework}

\subsection{Estancamiento de las interfaces de escritorio}
\label{subsec:stagnation}

\noindent
La interacción humano-computadora ha atravesado períodos de innovación genuina---la
metáfora de escritorio, la Web, la revolución táctil---que redefinieron los modelos
mentales de los usuarios \cite{Norman1988}. Sin embargo, en el escritorio las
últimas dos décadas constituyen una meseta funcional \cite{Jenson2024}.
Estudios recientes confirman que las herramientas integradas en los SO modernos
---exploradores de archivos, terminales, gestores de configuración---no han
experimentado cambios paradigmáticos desde mediados de los 2000
\cite{Grudin2017}. Para GNU/Linux, la CLI constituye la barrera principal
hacia una plataforma que domina la infraestructura profesional \cite{W3Techs2023}.

\subsection{Inteligencia Artificial Centrada en el Humano}
\label{subsec:hcai}

\noindent
Shneiderman \cite{Shneiderman2020, Shneiderman2022} formalizó HCAI como un
paradigma donde la IA \textit{amplifica} capacidades humanas en lugar de
reemplazarlas, articulado sobre cuatro principios: \textbf{transparencia} (el
usuario comprende las acciones del sistema), \textbf{control humano} (supervisión
y rechazo de propuestas), \textbf{seguridad} (prevención de acciones dañinas,
crítico donde un error puede causar pérdida irreversible de datos) y
\textbf{empoderamiento} (el usuario construye autonomía progresiva). Amershi et al.
\cite{Amershi2019} complementan con 18 directrices prácticas. HCAI provee un
cimiento coherente para repensar la interacción con el SO: un asistente fundado en
estos principios puede hacer accesible GNU/Linux asegurando que el usuario aprenda
en lugar de depender.

\subsection{Asistentes de voz y la brecha de interacción con el SO}
\label{subsec:voice-gap}

\noindent
Los avances en reconocimiento de voz---particularmente Whisper
\cite{Radford2022}---y en modelos de lenguaje de gran escala
\cite{Brown2020, OpenAI2023} han alcanzado madurez práctica multilingüe. NL2Bash
\cite{Lin2018} exploró la traducción de lenguaje natural a Bash (36\% de precisión
exacta) pero sin interacción en tiempo real ni validación de seguridad. Ningún
sistema existente integra reconocimiento de voz, ejecución de comandos,
clasificación de seguridad y retroalimentación educativa dentro de un flujo
orientado por HCAI. Esta brecha de integración define la oportunidad que aborda
este trabajo.


\section{Propuesta: BOTAS con Enfoque HCAI}
\label{sec:proposal}

\noindent
Los principios HCAI se materializan sobre GNU/Linux porque es el entorno donde
esta intervención genera mayor impacto: un sistema extraordinariamente potente cuya
principal barrera de adopción es la interacción, no la funcionalidad
\cite{Warschauer2004}. Su carácter abierto permite integración profunda, y su
gratuidad elimina barreras económicas---factor decisivo para contextos educativos
en América Latina.

BOTAS (\textit{Bot de Operaciones y Tareas Automatizadas del Sistema}) se originó
durante un programa de servicio social en una universidad pública mexicana,
explorando si el lenguaje natural podía servir como puente entre usuarios novatos
y la CLI de GNU/Linux. Desde un MVP dependiente de la nube, se incorporaron
iterativamente: detección de palabra de activación fuera de línea (Vosk),
reconocimiento local de intenciones, detección automática de distribución en
seis familias Linux y registro estructurado.

\subsection{Capacidades actuales}
\label{subsec:capabilities}

\noindent
BOTAS implementa un pipeline modular cuyo flujo se ilustra en la
Figura~\ref{fig:pipeline}: (1)~captura de entrada mediante GUI o palabra
de activación; (2)~detección local de intenciones por expresiones regulares;
(3)~análisis por GPT-4o para solicitudes complejas; (4)~validación de seguridad
determinista en tres niveles (\textit{seguro}, \textit{advertencia},
\textit{peligroso}); y (5)~ejecución con retroalimentación educativa y TTS.
Cada operación muestra el comando generado antes de ejecutarlo
(transparencia), explica su función (empoderamiento) y exige confirmación
para operaciones riesgosas (control humano).

Limitaciones actuales: dependencia de servicios en la nube (Whisper, GPT-4o),
dificultad con peticiones ambiguas y ausencia de diálogos multiturnos; se
desarrolla operación local con Vosk y modelos ligeros para entornos institucionales.

\begin{figure}[h]
\centering
\begin{tikzpicture}[
  node distance=4mm,
  box/.style={draw, rounded corners=2pt, minimum height=6mm, minimum width=22mm,
              font=\scriptsize, align=center, fill=#1},
  arr/.style={-{Stealth[length=1.5mm]}, thick},
]
\node[box=blue!10] (voz) {Voz / GUI};
\node[box=green!10, right=of voz] (local) {Intención\\local};
\node[box=orange!10, right=of local] (llm) {LLM};
\node[box=red!10, right=of llm] (seg) {Seguridad\\(3 niveles)};
\node[box=purple!8, right=of seg] (exec) {Ejecución\\+ TTS};
\draw[arr] (voz) -- (local);
\draw[arr] (local) -- node[above,font=\tiny]{compleja} (llm);
\draw[arr] (llm) -- (seg);
\draw[arr] (seg) -- (exec);
\draw[arr, dashed] (local) to[bend right=25] node[below,font=\tiny]{directa} (seg);
\end{tikzpicture}
\caption{Pipeline de procesamiento de BOTAS. Las peticiones rutinarias
(29\%) se resuelven localmente sin LLM (línea punteada).}
\label{fig:pipeline}
\end{figure}


\section{Evaluación Preliminar}
\label{sec:evaluation}

\subsection{Pruebas funcionales}
\label{subsec:testing}

\noindent
BOTAS fue evaluado con un conjunto de 200 tareas de administración de GNU/Linux
en cinco categorías: gestión de archivos~(45), información del sistema~(40),
control de procesos~(35), redes~(30) y permisos~(25). Las tareas fueron
seleccionadas a partir de operaciones frecuentes documentadas en manuales de
administración de sistemas y foros de soporte, cubriendo tres niveles de riesgo:
seguro~(169), advertencia~(18) y peligroso~(13). El conjunto fue compilado y
verificado manualmente por el desarrollador principal---restricción metodológica
que puede introducir sesgo de confirmación. La Tabla~\ref{tab:accuracy} resume
los resultados.

\begin{table}[h]
\caption{Precisión en la generación de comandos por categoría}
\label{tab:accuracy}
\centering
\begin{tabular}{lcc}
\toprule
\textbf{Categoría} & \textbf{Tareas} & \textbf{Precisión} \\
\midrule
Gestión de archivos    & 45 & 84\% \\
Información del sistema & 40 & 85\% \\
Control de procesos     & 35 & 83\% \\
Redes                   & 30 & 93\% \\
Permisos                & 25 & 80\% \\
\midrule
\textbf{Total}          & \textbf{200} & \textbf{89\%} \\
\bottomrule
\end{tabular}
\end{table}

Modos de error: fraseo ambiguo (6\%), errores de banderas en pipes múltiples
(3\%), indisponibilidad específica de distribución (2\%). La capa de seguridad
determinista detectó el 100\% de las 43 operaciones peligrosas (eliminación
recursiva, escalamiento de privilegios, modificaciones del sistema)---tasa
garantizada por expresiones regulares, independiente del modelo. El detector local
resolvió el 29\% de tareas sin LLM, con 100\% de precisión, ahorrando
500--1500\,ms por solicitud.

\subsection{Estudio piloto de usabilidad}
\label{subsec:usability}

\noindent
Tres usuarios sin experiencia en GNU/Linux evaluaron BOTAS en su etapa de producto
mínimo viable. Aunque insuficiente para generalización estadística, el piloto
aportó hallazgos cualitativos: los participantes valoraron la ejecución por
lenguaje natural, las explicaciones educativas y la retroalimentación de seguridad.
La retroalimentación constructiva identificó confusión con la sintaxis verbal,
frustración ante confirmaciones percibidas como excesivas y dificultad con
mensajes de error---informando directamente el rediseño hacia diálogos de
clarificación. Latencia: Whisper promedia 1--3\,s; Vosk responde en $<$200\,ms.


\section{Discusión y Conclusiones}
\label{sec:discussion}

\noindent
La precisión del 89\% implica que aproximadamente una de cada diez peticiones
genera un comando incorrecto. Sin embargo, el diseño HCAI mitiga el impacto
de este margen: cada comando se muestra al usuario antes de su ejecución
(transparencia), las operaciones riesgosas requieren confirmación explícita
(control humano) y el 100\% de operaciones destructivas son detectadas por la
capa de seguridad determinista. Así, un comando incorrecto no se ejecuta
automáticamente---el usuario conserva la decisión final, y el sistema
explica cada acción propuesta, convirtiendo el error en oportunidad educativa.
Desde la perspectiva latinoamericana, la barrera lingüística y los costos de
ecosistemas cerrados representan obstáculos para instituciones educativas
\cite{Warschauer2004}; la interacción por voz en español reduce estas barreras.

Trabajo futuro: consolidar la operación local, un estudio formal de usabilidad
con muestra ampliada que evalúe comprensión, percepción de control, confianza
y seguridad percibida, validación multidistribución y liberación como software
libre. Los principios HCAI materializados en BOTAS demuestran que las barreras
de acceso a sistemas potentes pueden reducirse sin sacrificar sus capacidades.


\section*{Agradecimientos}
\noindent
Este trabajo fue parcialmente desarrollado durante el Servicio Social en la
Universidad Tecnológica de la Mixteca. El autor agradece al Dr. Ricardo Ruiz
Rodríguez por su orientación y apoyo a lo largo del desarrollo de esta
investigación, así como a los participantes del estudio piloto de usabilidad
por su retroalimentación.

\bibliographystyle{sbc}
\bibliography{paper-es}

\end{document}
